\documentclass[10pt,a4paper]{amsart}
\usepackage[margin=19mm]{geometry}
\usepackage{amsmath,amssymb,mathtools}
\usepackage[scr=rsfs,cal=euler]{mathalfa}
\usepackage{xfrac}
\usepackage[unicode,hidelinks,bookmarks=false]{hyperref}
\newcommand{\ie}{i.e.\ }
\newcommand{\cf}{cf.\ }
\newcommand{\resp}{respectively\ }
\newcommand{\Subsection}{\subsection}
\providecommand{\nobreakdash}{\nobreak\hbox{-}}
\setlength{\emergencystretch}{2em}
\multlinegap0pt
\usepackage{bm}
\usepackage{bbm}
\usepackage{dsfont}
\usepackage{xspace}
\usepackage{cancel}
\usepackage{mathtools}
\usepackage{amsthm}
\makeatletter
\@ifundefined{theorem}{\newtheorem{theorem}{Theorem}}{}
\@ifundefined{lemma}{\newtheorem{lemma}[theorem]{Lemma}}{}
\@ifundefined{proposition}{\newtheorem{proposition}[theorem]{Proposition}}{}
\makeatother
\title[Discrete homotopy groups of cubical sets]{Discrete homotopy groups of cubical sets}
\author{Daisuke Kishimoto, Yichen Tong}
\date{Source: arXiv:2606.28693v3 (CC BY 4.0)}
\begin{document}
\maketitle

\noindent\emph{Source and licence.} Discrete homotopy groups of cubical sets. Version arXiv:2606.28693v3 (CC BY 4.0), distributed under the Creative Commons Attribution 4.0 International licence, as recorded in the arXiv licence metadata for this exact version and shown on the abstract page https://arxiv.org/abs/2606.28693v3. Full rights notice: licenses/arxiv-2606.28693-CC-BY-4.0.txt.\par\smallskip

\noindent\textit{Editorial scope.} Counted unit: the Proposition whose source label is \texttt{gluing grids} in the article (source0.tex lines 2403--2413, source label \texttt{gluing grids}), delivered together with the article's own complete original proof of it (source0.tex lines 2498--2624, one \texttt{proof} environment of 127 lines). No mathematics is written, repaired or re-derived: statement and proof are the article's own text line for line. The proof is detached from the statement in the source: the article states the result at lines 2403--2413 and prints its proof at lines 2498--2624, and the proof environment's own header names the statement, so the association is the article's own. Required context retained uncounted: common expansion (lines 2456--2459); stabilization (lines 2471--2478). Every definition, display and label of those lines is retained unchanged. Omitted from this package and not redistributed: 1--2402, 2414--2455, 2460--2470, 2479--2497, 2625--3597. Nothing inside a retained excerpt is omitted or altered: every retained body is delivered exactly as the article's lines are written, so no omission receipt of any kind accompanies this package. Citations: the retained excerpts carry no citation command at all, so no bibliography entry of the article is transcribed and no reference is redistributed. Reference adaptation: none was needed and none was made; every reference inside the delivered unit resolves inside it, so no unresolved reference or citation is produced. Printed slips kept literal: no printed word, symbol, hypothesis or number of the counted statement or of its proof was altered. Layout and class: the article's own class is \texttt{amsart} with packages including amsmath, amssymb, amsthm, amsfonts, newtxtext, newtxmath, mathalfa, xy, geometry, tikz, hyperref; the wrapper is the lane's pinned \texttt{amsart} editorial wrapper at 10pt on A4, which restates the theorem-like environments the retained text needs and adds the article's own macro definitions that the retained text needs (none), with extra wrapper packages (bm, bbm, dsfont, xspace, cancel, mathtools). The delivered numbering is the unit's own. Assets: no retained span contains a figure environment, a graphics-inclusion command, a PDF-page inclusion, a table or a TikZ picture; no image file is copied into this package, the package declares no asset dependency and no image file is redistributed with it. Licence: version arXiv:2606.28693v3 of this article is distributed under the Creative Commons Attribution 4.0 International licence, as recorded in the arXiv licence metadata for this exact version and shown on the abstract page https://arxiv.org/abs/2606.28693v3. A full rights notice is delivered at licenses/arxiv-2606.28693-CC-BY-4.0.txt. Characters: every retained excerpt is pure ASCII, so the delivered engine, XeLaTeX with its default Latin Modern text font, reports no missing character.
	\begin{lemma}
		\label{common expansion}
		Let $X$ be a cubical set. For $x,y\in\mathcal{G}X_n$ such that $x\simeq_0y$, there are $n$-expanders $\epsilon$ and $\delta$ such that $x^\epsilon=y^\delta$.
	\end{lemma}

	\begin{lemma}
		\label{stabilization}
		Let $x\in\mathcal{G}X_n$, and let $\epsilon$ and $\epsilon'$ be $n$-expanders. If $x^{\epsilon}=x^{\epsilon'}$, then there exist $n$-expanders $\delta,\delta'$ such that
		\[
		\epsilon\circ\delta=\epsilon'\circ\delta'.
		\]
		Moreover, if $\epsilon$ and $\epsilon'$ are supported within $S\subset\{1,\ldots,n\}$, then $\delta$ and $\delta'$ may be chosen to be supported within $S$.
	\end{lemma}

	\begin{proposition}
		\label{gluing grids}
		There exists a family $\{v_{p_1,\ldots,p_n}\}_{p_1,\ldots,p_n\in\{-1,0,1\}}$ in $\mathcal{G}X_{n+1}$ such that for all $p_1,\ldots,p_n$,
		\[
		v_{p_1,\ldots,p_n}\simeq_0 u_{p_1,\ldots,p_n}
		\]
		and
		\[
		v_{p_1,\ldots,p_{i-1},0,p_{i+1},\ldots,p_n}\partial_{i,\alpha}=v_{p_1,\ldots,p_{i-1},(-1)^{1-\alpha},p_{i+1},\ldots,p_n}\partial_{i,1-\alpha}.
		\]
	\end{proposition}

	\begin{proof}
		[Proof of Proposition \ref{gluing grids}]
		We aim to construct $(n+1)$-expanders $\epsilon_{p_1,\ldots,p_n}$ for all $p_1,\ldots,p_n$ such that
		\[
		v_{p_1,\ldots,p_n}=(u_{p_1,\ldots,p_n})^{\epsilon_{p_1,\ldots,p_n}}
		\]
		where $v_{p_1,\ldots,p_n}$ satisfy the condition in the statement. The construction is obtained by applying Lemma \ref{common expansion} iteratively such that $v_{p_1,\ldots,p_n}$ is sufficiently refined in every coordinate and compatible with $v_{q_1,\ldots,q_n}$, where $|p_i-q_i|\le1$ for $i=1,...,n$. For $0\le k\le n$, let
		\[
		\mathcal{S}(k,n)=\{(s_1,\ldots,s_k)\in\{1,\ldots,n\}^k\mid s_i\ne s_j\text{ for }i\ne j\}.
		\]
		For $S\in\mathcal{S}(k,n)$, let $\underline{S}=\{j_1<\cdots<j_k\}$ denote the underlying poset of $S$. We construct a series of $(n+1)$-expanders $\epsilon_{p_1,\ldots,p_n}^S$ for all $p_1,\ldots,p_n$ and all $S\in\mathcal{S}(k,n)$ with $0\le k\le n$. Define $\epsilon_{p_1,\ldots,p_n}^\emptyset$ as the identity expander. Suppose that we have constructed $(n+1)$-expanders $\epsilon_{p_1,\ldots,p_n}^S$ for all $S\in\bigsqcup_{l=0}^{k-1}\mathcal{S}(l,n)$. Set
		\[
		u_{p_1,\ldots,p_n}^S=(u_{p_1,\ldots,p_n})^{\epsilon_{p_1,\ldots,p_n}^S}.
		\]
		Let $S=(s_1,\ldots,s_k)\in\mathcal{S}(k,n)$ with $\underline{S}=\{j_1<\cdots<j_k\}$. Set $\widehat{S}=(s_1,\ldots,s_{k-1})\in\mathcal{S}(k-1,n)$. Fix $p_1,\ldots,p_n\in\{-1,0,1\}$ such that
		\[
		\{j\in\{1,\ldots,n\}\mid p_j=0\}=\underline{S}.
		\]
		%		and let $\{i_1<\cdots<i_{n-k+1}\}=\{1,\ldots,n+1\}-\{j_1<\cdots<j_k\}$.
		For $p=\pm 1$, let
		\[
		\alpha(p)=
		\begin{cases}
			0&(p=-1)\\
			1&(p=1).
		\end{cases}
		\]
		Choose $p_{j_1}',\ldots,p_{j_k}'\in\{-1,1\}$. Set
		\begin{align*}
			&N(u_{p_1,\ldots,p_n}\partial_{j_k,\alpha(p_{j_k}')}\cdots\partial_{j_1,\alpha(p_{j_1}')})\\
			&=\{u_{q_1,\ldots,q_n}\mid q_{i_l}=p_{i_l}\;(l=1,\ldots,n-k+1)\text{ and }q_{j_m}\in\{0,p_{j_m}'\}\;(m=1,\ldots,k)\}.
		\end{align*}
		Then
		\[
		u_{q_1,\ldots,q_n}^T\partial_{j_k,\alpha(q_{j_k})}\cdots\partial_{j_1,\alpha(q_{j_1})}\simeq_0u_{p_1,\ldots,p_n}\partial_{j_k,\alpha(p_{j_k}')}\cdots\partial_{j_1,\alpha(p_{j_1}')}
		\]
		for all $q_1,\ldots,q_n$ and all $T\in\bigsqcup_{l=0}^{k-1}\mathcal{S}(l,n)$. By Lemma \ref{common expansion}, there exists $w_{p_1,\ldots,p_n}^{p_{j_1}',\ldots,p_{j_k}'}\in\mathcal{G}X_{n-k+1}$ such that
		\[
		w_{p_1,\ldots,p_n}^{p_{j_1}',\ldots,p_{j_k}'}\ge u_{q_1,\ldots,q_n}^T\partial_{j_k,\alpha(q_{j_k})}\cdots\partial_{j_1,\alpha(q_{j_1})}
		\]
		for all $u_{q_1,\ldots,q_n}\in N(u_{p_1,\ldots,p_n}\partial_{j_k,\alpha(p_{j_k}')}\cdots\partial_{j_1,\alpha(p_{j_1}')})$ and all $T\in\bigsqcup_{l=0}^{k-1}\mathcal{S}(l,n)$. Since
		\[
		u_{p_1,\ldots,p_n}\in\bigcap_{p_{j_1}',\ldots,p_{j_k}'\in\{-1,1\}}N(u_{p_1,\ldots,p_n}\partial_{j_k,\alpha(p_{j_k}')}\cdots\partial_{j_1,\alpha(p_{j_1}')}),
		\]
		we may assume
		\[
		w_{p_1,\ldots,p_n}^{p_{j_1}',\ldots,p_{j_k}'}=(u_{p_1,\ldots,p_n}\partial_{j_k,\alpha(p_{j_k}')}\cdots\partial_{j_1,\alpha(p_{j_1}')})^\epsilon
		\]
		for some $(n-k+1)$-expander $\epsilon$ that does not depend on the choice of $p_1',\ldots,p_k'$. Fix $p_{j_1}',\ldots,p_{j_k}'\in\{-1,1\}$ and $u_{q_1,\ldots,q_n}\in N(u_{p_1,\ldots,p_n}\partial_{j_k,\alpha(p_{j_k}')}\cdots\partial_{j_1,\alpha(p_{j_1}')})$. Then there exists an $(n+1)$-expander $\epsilon_{q_1,\ldots,q_n}^{\widehat{S},S}$ supported within $\underline{S}$, which does not depend on the choice of $p_{j_1}',\ldots,p_{j_k}'$, such that
		\[
		(u_{q_1,\ldots,q_n}^{\widehat{S}})^{\epsilon_{q_1,\ldots,q_n}^{\widehat{S},S}}\partial_{j_k,\alpha(p_{j_k}')}\cdots\partial_{j_1,\alpha(p_{j_1}')}=w_{p_1,\ldots,p_n}^{p_{j_1}',\ldots,p_{j_k}'}.
		\]
		We now define
		\[
		\epsilon_{q_1,\ldots,q_n}^{S}=\epsilon_{q_1,\ldots,q_n}^{\widehat{S}}\circ\epsilon_{q_1,\ldots,q_n}^{\widehat{S},S}.
		\]
		%Note that we may choose the expanders appropriately such that it does not depend on $p_{j_1},\ldots,p_{j_k}\in\{-1,1\}$.
		Then for any elements $u_{q_1,\ldots,q_n}$ and $u_{q_1',\ldots,q_n'}$ of $N(u_{p_1,\ldots,p_n}\partial_{j_k,\alpha(p_{j_k}')}\cdots\partial_{j_1,\alpha(p_{j_1}')})$ and all $T\in\bigsqcup_{l=0}^{k}\mathcal{S}(l,n)$, we have
		\[
		u_{q_1,\ldots,q_n}^S\partial_{j_k,\alpha(q_{j_k})}\cdots\partial_{j_1,\alpha(q_{j_1})}\ge u_{q_1',\ldots,q_n'}^T\partial_{j_k,\alpha(q_{j_k}')}\cdots\partial_{j_1,\alpha(q_{j_1}')}.
		\]
		In particular, if $T\in\mathcal{S}(k,n)$ satisfies $\underline{S}=\underline{T}$, then
		\begin{equation}
			\label{useful}
			u_{q_1,\ldots,q_n}^S\partial_{j_k,\alpha(q_{j_k})}\cdots\partial_{j_1,\alpha(q_{j_1})}=u_{q_1',\ldots,q_n'}^T\partial_{j_k,\alpha(q_{j_k}')}\cdots\partial_{j_1,\alpha(q_{j_1}')}.
		\end{equation}
		Let $q_1=q_1',\ldots,q_n=q_n'$ in \eqref{useful}. By Lemma \ref{stabilization}, we may choose $\epsilon_{q_1,\ldots,q_n}^S$ for $S\in\mathcal{S}(k,n)$ such that
		\[
		(\epsilon_{q_1,\ldots,q_n}^S)_i=(\epsilon_{q_1,\ldots,q_n}^T)_i
		\]
		for all $i\in\{1,\ldots,n+1\}-\underline{S}$ and all $T\in\mathcal{S}(k,n)$ with $\underline{S}=\underline{T}$, where
		$$\epsilon_{q_1,\ldots,q_n}^S=((\epsilon_{q_1,\ldots,q_n}^S)_1,\ldots,(\epsilon_{q_1,\ldots,q_n}^S)_{n+1}).$$ 
		Therefore, $(\epsilon_{q_1,\ldots,q_n}^S)_i$ for $i\in\{1,\ldots,n+1\}-\{j_1,\ldots,j_k\}$ depends only on $\underline{S}$.
		
		
		
		Define an $(n+1)$-expander $\epsilon_{q_1,\ldots,q_n}$ by
		\[
		(\epsilon_{q_1,\ldots,q_n})_i=
		\begin{cases}
			(\epsilon_{q_1,\ldots,q_n}^{(1,\ldots,i-1,i+1,\ldots,n)})_i&(i=1,\ldots,n)\\
			(\epsilon_{q_1,\ldots,q_n}^{(1,\ldots,n)})_{n+1}&(i=n+1)
		\end{cases}
		\]
		and set
		\[
		v_{p_1,\ldots,p_n}=(u_{p_1,\ldots,p_n})^{\epsilon_{p_1,\ldots,p_n}}.
		\]
		It remains to show that
		\begin{equation}
			\label{goal}
			v_{p_1,\ldots,p_{i-1},0,p_{i+1},\ldots,p_n}\partial_{i,\alpha}=v_{p_1,\ldots,p_{i-1},(-1)^{1-\alpha},p_{i+1},\ldots,p_n}\partial_{i,1-\alpha}.
		\end{equation}
		Let $T=(i)\in\mathcal{S}(1,n)$. Then
		\begin{equation}
			\label{useful2}
			u_{p_1,\ldots,p_{i-1},0,p_{i+1},\ldots,p_n}^T\partial_{i,\alpha}=u_{p_1,\ldots,p_{i-1},(-1)^{1-\alpha},p_{i+1},\ldots,p_n}^T\partial_{i,1-\alpha}.
		\end{equation}
		For $j=1,\ldots,n+1$, define $S_j=(s_1,\ldots,s_k)\in\mathcal{S}(k,n)$ with $k=n-1$ or $n$ by $s_1=i$ and
		\[
		\underline{S_j}=
		\begin{cases}
			\{1<\cdots<j-1<j+1<\cdots<n\}&(j=1,\ldots,n)\\
			\{1<\cdots<n\}&(j=n+1)
		\end{cases}
		\]
		Then we have
		\[
		(\epsilon_{q_1,\ldots,q_n})_j=(\epsilon_{q_1,\ldots,q_n}^{(i)}\circ\epsilon_{q_1,\ldots,q_n}^{(i),S_j})_j.
		\]
		Hence, it suffices to show that
		\begin{equation}
			\label{goal2}
			(\epsilon_{p_1,\ldots,p_{i-1},0,p_{i+1},\ldots,p_n}^{(i),S_j})_j=(\epsilon_{p_1,\ldots,p_{i-1},\pm1,p_{i+1},\ldots,p_n}^{(i),S_j})_j
		\end{equation}
		for $j\ne i$. We only prove the case that $j\le n$ because the case that $j=n+1$ is identical. By \eqref{useful},
		\begin{align*}
			&u_{p_1,\ldots,p_{i-1},0,p_{i+1},\ldots,p_n}^{S_j}\partial_{n,\alpha(p_n)}\cdots\partial_{j+1,\alpha(p_{j+1})}\partial_{j-1,\alpha(p_{j-1})}\cdots\partial_{1,\alpha(p_1)}\\
			&=u_{p_1,\ldots,p_{i-1},(-1)^{1-\alpha},p_{i+1},\ldots,p_n}^{S_j}\partial_{n,\alpha(p_n)}\cdots\partial_{j+1,\alpha(p_{j+1})}\partial_{j-1,\alpha(p_{j-1})}\cdots\partial_{1,\alpha(p_1)}
		\end{align*}
		and by \eqref{useful2},
		\begin{align*}
			&u_{p_1,\ldots,p_{i-1},0,p_{i+1},\ldots,p_n}^{(i)}\partial_{n,\alpha(p_n)}\cdots\partial_{j+1,\alpha(p_{j+1})}\partial_{j-1,\alpha(p_{j-1})}\cdots\partial_{1,\alpha(p_1)}\\
			&=u_{p_1,\ldots,p_{i-1},(-1)^{1-\alpha},p_{i+1},\ldots,p_n}^{(i)}\partial_{n,\alpha(p_n)}\cdots\partial_{j+1,\alpha(p_{j+1})}\partial_{j-1,\alpha(p_{j-1})}\cdots\partial_{1,\alpha(p_1)}.
		\end{align*}
		Hence by Lemma \ref{stabilization}, \eqref{goal} follows.
	\end{proof}

\end{document}
